From \parencite{chen2023ml_manufacturing_industry4} :

"Despite numerous ML studies and their promising performance, it remains very difficult for non-experts working in the manufacturing industry to begin developing ML solutions for their specific problems."


"The first challenging part of application is to formulate the actual problems to be solved."


"In spite of the enormous number of ML use cases, there is no guidance or standard for developing ML solutions from ideation to deployment."


"The selection of an ML model plays a critical role in the outcome."


"To fill the gap between academic research and manufacturing industries, we provided an actionable pipeline for the implementation of ML solutions by production engineers from ideation through to deployment."


"We hope that this can provide support in method selection for decision-makers considering ML solutions."

From \parencite{sala2018selecting_ml_algorithm} :

"The need of synthesis and guidelines to drive the selection of the most suitable algorithm for a specific scope arises."

Om begrensninger ved eksisterende tilnærminger (f.eks. Azure):

"Some of the frameworks currently available tend to use as main selection drivers the accuracy and/or prediction speed of the ML algorithms. The problem in this class of selection frameworks is that they may result too general most of the time."

"Most of the times selecting an algorithm only on the base of the promised accuracy or computational speed leads to unsatisfactory results."

Om begrensninger ved statistiske kriterier (AIC, BIC):

"The problem with these criteria is that they require that the user knows the characteristics of each model and that s/he computes the value for each single model under consideration. In this case, the model selection could be biased by the initial pool of algorithms selected by the user."

"A unique ML algorithm able to deal with the analysis of every dataset does not exist, making the selection of the right ML algorithm challenging."

"The presence in literature of a wide range of different algorithms can create disorientation in the users not acquainted with their known strengths and weaknesses."

"This selection framework aims at supporting the researchers who are approaching the ML field guiding them through the selection of a set of algorithms suitable for their scopes, helping them to avoid incurring in the application of algorithms which are not suitable for the data they are dealing with."

"The proposed selection framework does not claim to be complete with respect to the available literature on the topic and/or to suggest the best algorithm to the user due to the vast complexity characterising the ML field of research and application."

"In the selection of a suitable ML algorithm for data analysis many aspects should be taken into account, and most of the times selecting an algorithm only on the base of the promised accuracy or computational speed leads to unsatisfactory results."

From \parencite{arena2024conceptual_framework_ml_algorithm_selection_predictive_maintenance} :

"The need to choose the right algorithm for a specific task arises as a challenging issue since it is considered an essential stage in the development and implementation of an ML-oriented approach."

"There is no unique ML algorithm that can efficiently handle the analysis of every dataset or all use cases."

"Often, the choice is established merely by the comparison between different algorithms whose final decision falls on the one that ensures the 'best solution' only in terms of predictive and computational performance. As a consequence, this may frequently lead to results that are poor or too generic for the problem at hand."

"It is crucial to define the applicability domain in terms of technical, organizational, and economic assessment aiming at assisting users during the decision-making process."

"Most of the articles focus on the development of the ML model, based on specific cases and applications. Most of them propose ML approaches with a predefined maintenance objective, not considering what conditions lead to selecting one rather than another."

"Many of them focus on accuracy, prediction speed, or computational requirements."

"No tool relates algorithmic approaches to the maintenance process."

"The theoretical accuracy and the computational speed do not ensure satisfactory results, making the algorithm selection a non-trivial task."

"The basic idea is to develop an easy and responsive tool aiming at performing agile and informed decision-making in establishing a proper ML algorithm suitable for the specific case at hand."

"Implementing an ML-based PdM strategy presents multiple challenges, such as the investment costs, the connection of physical assets, the storage of the data, extraction of valuable information, and the development of precise ML algorithms."

From \parencite{wolpert1997no_free_lunch} :

For any algorithm, any elevated performance over one class of problems is offset by performance over another class."

"If some algorithm's performance is superior to that of another algorithm over some set of optimization problems, then the reverse must be true over the set of all other optimization problems."

"If knowledge of problem characteristics is not incorporated into the optimization algorithm, then there are no formal assurances that the algorithm chosen will be at all effective."

From \parencite{smith_miles2009algorithm_selection} :

"The Algorithm Selection Problem seeks to answer the question: which algorithm is likely to perform best for my problem?"

"From the early work of John Rice which formalized the Algorithm Selection Problem, to the No Free Lunch Theorems of Wolpert and Macready which stated that for any algorithm, any elevated performance over one class of problems is exactly paid for in performance over another class, it has long been acknowledged that there is no universally-best algorithm to tackle a broad problem domain."

"The NFL theorems suggest that we need to move away from a 'black-box' approach to algorithm selection. We need to understand more about the characteristics of the problem in order to select the most appropriate algorithm."

From 